\section{Introduction}
\label{sec:intro}

Scientific discovery increasingly occurs at disciplinary boundaries.
The mechanisms by which a gene mutation in cancer signaling alters metabolic chemistry,
or how a neurotransmitter synthesis enzyme predicts pharmacological drug sensitivity,
are precisely the cross-domain questions that no single-domain knowledge graph (KG)
can answer in isolation.

Existing approaches to KG-based discovery---link prediction~\cite{TODO-TransE,TODO-RotatE},
KG completion~\cite{TODO-ComplEx}, and semantic similarity methods---operate within
a single KG schema.
They can predict missing edges within a domain but cannot generate hypotheses that
cross the boundary between two independently curated domain graphs with disjoint
node vocabularies.

We address this gap with a multi-operator pipeline that treats cross-domain discovery
as a graph composition problem.
The key operator is \textbf{align}: it identifies bridge entities shared between two
domain KGs (e.g., a metabolite that appears in both a biological signaling graph and
a chemical reaction graph under different identifiers), collapses them into shared
nodes, and enables the \textbf{compose} operator to enumerate multi-hop paths that
cross the domain boundary.
Without alignment, compose produces zero cross-domain candidates; with alignment,
it produces structurally unique paths connecting biological genes and proteins
to chemical reaction classes and functional groups.

The pipeline builds on Swanson's ABC discovery model~\cite{swanson1986fish}, which
finds implicit connections $A \to C$ via a shared concept $B$ in two literatures.
Our contribution is to formalise this as a graph-algebraic operator pipeline,
to study its empirical behaviour on real Wikidata-derived knowledge graphs, and to
characterise the conditions under which its output is scientifically meaningful.

\subsection*{Contributions}

This paper makes three empirical contributions:

\begin{enumerate}[leftmargin=*, label=\textbf{\arabic*.}]

\item \textbf{Alignment is necessary for cross-domain reachability.}
Without bridge alignment, the compose operator produces zero cross-domain
candidate pairs.
With alignment, it produces 5--55 alignment-dependent unique cross-domain paths
per domain pair, scaling non-linearly with KG size.
Replicated across 3 independent bio-chem domain pairs.

\item \textbf{Quality filtering makes raw deep output useful.}
Raw 3-hop+ cross-domain candidates from real KGs contain
25\% semantically drifted content (structural-chemical expansions
rather than mechanistic hypotheses).
A relation-type filter derived from qualitative review of one domain pair
eliminates all drift while preserving all promising candidates, and
transfers across two independent domain pairs without parameter adjustment.

\item \textbf{Bridge structural diversity, not bridge count, drives candidate yield.}
Across three domain pairs with 5--9 aligned bridge pairs,
unique-to-alignment candidate counts vary from 5 to 55.
This variation is explained by bridge entity dispersion (the number of
distinct downstream nodes each bridge class connects to), not by bridge count.

\end{enumerate}

We also contribute a negative finding: a post-hoc relation semantics audit
reveals that all five qualitatively reviewed candidates carry at least one
relation edge where the class-level KG encoding cannot validate
substrate-specific chemical or biological claims.
This finding unifies previously separate candidate weaknesses under a single
root cause and is the primary open problem for future work.

The remainder of this paper is structured as follows.
Section~\ref{sec:related} reviews related work.
Section~\ref{sec:method} describes the pipeline operators and scoring rubric.
Section~\ref{sec:setup} presents the experimental setup and data.
Section~\ref{sec:results} reports the main results for Claims 1--3.
Section~\ref{sec:casestudies} presents five representative candidate case studies.
Section~\ref{sec:discussion} interprets the findings, with relation semantic
insufficiency as the central theme.
Sections~\ref{sec:threats}--\ref{sec:limitations} address threats to validity
and limitations.
Section~\ref{sec:conclusion} concludes with future work priorities.
